% Source: physical PDF pp. 323--327 (printed pp. 300--304). Coverage: complete Section 21.2, from its heading through immediately before Section 21.3; Eqs. (21.2.1)--(21.2.23), citations [4], [5], and [5a], no figures, tables, or footnotes. Uncertainty: none; source anomalies are recorded in the wave report.
\subsection{\texorpdfstring{Renormalizable $\xi$-Gauges}{Renormalizable xi-Gauges}}
\label{sec:21.2}

In 1971 't Hooft~\hyperref[ch21-ref-4]{[4]} showed that path integrals in
spontaneously broken gauge theories can be calculated in a gauge in which
the vector boson propagators vanish for momentum $k\to\infty$ as $k^{-2}$,
so that these theories satisfy the power-counting test for renormalizability
described in Section 12.2. Here we shall describe a larger class of
renormalizable gauges, parameterized by an arbitrary constant $\xi$, that
was introduced a little later by Fujikawa, Lee, and
Sanda.~\hyperref[ch21-ref-5]{[5]}

In general gauges the kinematic term \eqref{eq:21.1.3} in the Lagrangian for
the scalar fields of a theory contains a cross term
\begin{equation*}
 i\sum_{nma}\partial_\mu\phi'_n
 t^a_{nm}A_a^\mu v_m,
\end{equation*}
where $v_m$ is the vacuum expectation value of $\phi_m$, and $\phi'_n$ is
the shifted field defined by Eq.~\eqref{eq:21.1.4}. In unitarity gauge this
term vanished as a consequence of the gauge condition \eqref{eq:21.1.2}.
We will adopt a different approach here, similar to that of Sections 15.5
and 15.6. A functional $B[f]$ is introduced into the path integral, with
\begin{equation}
 B[f]=\exp\left(
 -\frac{i}{2\xi}\int d^4x\sum_a f_a f_a
 \right).
 \tag{21.2.1}\label{eq:21.2.1}
\end{equation}
This is equivalent to adding a gauge-fixing term in the Lagrangian
\begin{equation}
 \mathcal L_{\rm gf}=-\frac1{2\xi}f_a f_a.
 \tag{21.2.2}\label{eq:21.2.2}
\end{equation}
Instead of taking $f_a=\partial_\mu A_a^\mu$ as in Section 15.5,
we shall now take the gauge-fixing function as
\begin{equation}
 f_a=\partial_\mu A_a^\mu
 -i\xi(t_a)_{nm}\phi'_n v_m,
 \tag{21.2.3}\label{eq:21.2.3}
\end{equation}
which is designed so that the above cross term in \eqref{eq:21.1.3} is
cancelled by the cross term in Eq.~\eqref{eq:21.2.2}. Unitarity gauge is now
a special case; for $\xi\to\infty$, the gauge-fixing functional
\eqref{eq:21.2.1} is infinitely sharply peaked at a $\phi'$ that satisfies
the unitarity gauge condition \eqref{eq:21.1.2}. Another special case is
provided by the limit $\xi\to0$; here the gauge-fixing functional is peaked
at a gauge field that satisfies the Landau gauge condition
$\partial_\mu A_a^\mu=0$.

We also include in the Lagrangian a quartic polynomial $-P(\phi)$, subject
to the gauge-invariance condition
\begin{equation}
 \frac{\partial P(\phi)}{\partial\phi_n}
 (t_a)_{nm}\phi_m=0.
 \tag{21.2.4}\label{eq:21.2.4}
\end{equation}
We must of course also include in the Lagrangian a gauge-field term
\begin{equation}
 \mathcal L_A=-\frac14\sum_a
 F_a^{\mu\nu}F_{a,\mu\nu}.
 \tag{21.2.5}\label{eq:21.2.5}
\end{equation}
The total Lagrangian density of gauge fields and scalars is then
\begin{align}
 \mathcal L_{A,\phi}
 \equiv{}&\mathcal L_A+\mathcal L_\phi+\mathcal L_{\rm gf}
 \notag\\
 ={}&-\frac14\sum_a F_a^{\mu\nu}F_{a,\mu\nu}
 +\frac12\sum_{nab}
 (t^a\phi)_n(t^b\phi)_n
 A_a^\mu A_{b\mu}
 \notag\\
 &-\frac1{2\xi}\sum_a
 (\partial_\mu A_a^\mu)(\partial_\nu A_a^\nu)
 -\frac12\sum_n\partial_\mu\phi'_n\partial^\mu\phi'_n
 \notag\\
 &+\frac{\xi}{2}\sum_{a nm}
 (t_a v)_n(t_a v)_m\phi'_n\phi'_m
 -P(\phi)
 \notag\\
 &+i\sum_{a nm}
 \partial_\mu\phi'_n(t_a)_{nm}\phi'_mA_a^\mu
 +\text{total derivatives}.
 \tag{21.2.6}\label{eq:21.2.6}
\end{align}

As we saw in Section 15.6, the introduction of a gauge-fixing functional
$B[f]$ requires also the introduction of a ghost field $\omega_a(x)$,
with Lagrangian depending on the gauge-transformation properties of
$f_a$. Under a general gauge transformation (with an arbitrary
function $\epsilon_a(x)$)
\begin{equation}
 \delta A_a^\mu
 =-\sum_{bc}C_{abc}
 \epsilon_b A_c^\mu+\partial^\mu\epsilon_a,
 \tag{21.2.7}\label{eq:21.2.7}
\end{equation}
\begin{equation}
 \delta\phi_n=i\sum_a
 \epsilon_a(t_a)_{nm}\phi_m,
 \tag{21.2.8}\label{eq:21.2.8}
\end{equation}
we have
\begin{align}
 \delta f_a={}&\Box\epsilon_a
 -\sum_{bc}C_{abc}
 \partial_\mu(\epsilon_b A_c^\mu)
 \notag\\
 &+\xi\sum_{nb}(t_a v)_n
 \epsilon_b(t_b\phi)_n.
 \tag{21.2.9}\label{eq:21.2.9}
\end{align}
According to the general results of Section 15.6, this yields a ghost
Lagrangian
\begin{align}
 \mathcal L_\omega=\omega_a^*\Bigg[
 &\Box\omega_a
 -\sum_{bc}C_{abc}
 \partial_\mu(\omega_b A_c^\mu)
 \notag\\
 &+\xi\sum_{nb}(t_a v)_n
 \omega_b(t_b\phi)_n\Bigg].
 \tag{21.2.10}\label{eq:21.2.10}
\end{align}
Finally, if the theory involves spin $\frac12$ fermions there will also be a
general renormalizable term
\begin{equation}
 \mathcal L_\psi
 =-\bar\psi\left(
 \sl{\partial}-i\sl{A}_a t_a^{(\psi)}
 +m_0+\Gamma_n\phi_n\right)\psi,
 \tag{21.2.11}\label{eq:21.2.11}
\end{equation}
where $t_a^{(\psi)}$ is the matrix representation of the generators of
the gauge group for the fermions (including coupling-constant factors), and
$m$ and $\Gamma_n$ are constant matrices (in general, linear combinations
of terms proportional to Dirac matrices $\mathbf{1}$ and $\gamma_5$)
satisfying the gauge-invariance conditions
\begin{equation}
 [t_a^{(\psi)},\gamma_4m_0]=0,
 \tag{21.2.12}\label{eq:21.2.12}
\end{equation}
\begin{equation}
 [t_a^{(\psi)},\gamma_4\Gamma_n]
 +\sum_m(t_a)_{mn}\gamma_4\Gamma_m=0.
 \tag{21.2.13}\label{eq:21.2.13}
\end{equation}
(The factor $\gamma_4\equiv i\gamma^0$ arises from the definition
$\bar\psi\equiv\psi^\dagger\gamma_4$. It is relevant only if
$t_a^{(\psi)}$ involves terms proportional to $\gamma_5$.) The general
theorem proved in Sections 15.5 and 15.6 guarantees that the $S$-matrix
calculated from a Lagrangian given by the sum of Eqs.~\eqref{eq:21.2.6},
\eqref{eq:21.2.10}, and \eqref{eq:21.2.11} is independent of the choice of
the parameter $\xi$ appearing in the gauge-fixing function
\eqref{eq:21.2.3}, and so for any $\xi$ will give the same results as the
choice $\xi=\infty$ corresponding to unitarity gauge.

To derive the propagators for all these fields, we need the part of the
Lagrangian that is quadratic in fields:
\begin{align}
 \mathcal L_{\mathrm{QUAD}}={}&
 -\frac14\sum_a
 (\partial^\mu A_a^\nu-\partial^\nu A_a^\mu)
 (\partial_\mu A_{a\nu}-\partial_\nu A_{a\mu})
 \notag\\
 &-\frac12\sum_{ab}\mu^2_{ab}
 A_a^\mu A_{b\mu}
 -\frac1{2\xi}\sum_a
 (\partial_\mu A_a^\mu)(\partial_\nu A_a^\nu)
 \notag\\
 &-\frac12\sum_n(\partial_\mu\phi'_n)(\partial^\mu\phi'_n)
 -\frac12\sum_{nm}M^2_{nm}\phi'_n\phi'_m
 -\bar\psi(\sl{\partial}+m)\psi
 \notag\\
 &-\partial_\mu\omega_a^*\partial^\mu\omega_a
 -\xi\sum_{ab}\mu^2_{ab}
 \omega_a^*\omega_b
 +\text{total derivatives},
 \tag{21.2.14}\label{eq:21.2.14}
\end{align}
where $\mu^2_{ab}$ is the vector boson mass matrix
\eqref{eq:21.1.7}:
\begin{equation}
 \mu^2_{ab}
 =-\sum_n(t^a v)_n(t^b v)_n,
 \tag{21.2.15}\label{eq:21.2.15}
\end{equation}
and $M^2_{nm}$ and $m$ are new scalar and fermion mass matrices:
\begin{align}
 M^2_{nm}
 ={}&\left.
 \frac{\partial^2P(\phi)}{\partial\phi_n\partial\phi_m}
 \right|_{\phi=v}
 -\frac{\xi}{2}\sum_a
 (t_a v)_n(t_a v)_m,
 \notag\\
 m={}&m_0+\sum_n\Gamma_n v_n.
 \tag{21.2.16}\label{eq:21.2.16}
\end{align}
We see from Eq.~\eqref{eq:21.2.14} that the ghosts have gauge-dependent
masses, equal to $\sqrt{\xi}$ times the corresponding vector boson masses.

These expressions give the particle masses in the zeroth order of
perturbation theory. To this order, the vacuum expectation value $v_n$ is
just the location of the minimum of the polynomial `potential' $P(\phi)$:
\begin{equation}
 \left.\frac{\partial P(\phi)}{\partial\phi_n}\right|_{\phi=v}=0.
 \tag{21.2.17}\label{eq:21.2.17}
\end{equation}
Also, as we saw in Section 19.2, it follows from Eqs.~\eqref{eq:21.2.4} and
\eqref{eq:21.2.17} that
\begin{equation}
 \sum_m\left.
 \frac{\partial^2P(\phi)}{\partial\phi_n\partial\phi_m}
 \right|_{\phi=v}(t_a v)_m=0
 \tag{21.2.18}\label{eq:21.2.18}
\end{equation}
for all $a$. It follows then that in place of the Goldstone modes with
mass zero, the scalar boson mass-squared matrix in Eq.~\eqref{eq:21.2.16}
has eigenvalues equal to $\sqrt{\xi}$ times the non-zero vector boson
masses. That is, if $\mu^2_{ab}$ ha an eigenvector $c_b$ with
eigenvalue $\mu^2$, then $\sum_b c_b t_b v$ is an eigenvector
of $M^2$ with an eigenvalue $\xi\mu^2$:
\begin{align}
 \sum_mM^2_{nm}
 \left(\sum_b c_b t_b v\right)_m
 &=
 \xi\sum_{ab}\mu^2_{ab}c_b
 (t_a v)_n
 \notag\\
 &=\xi\mu^2
 \left(\sum_a c_a t_a v\right)_n.
 \tag{21.2.19}\label{eq:21.2.19}
\end{align}
The other eigenvectors of $M^2_{nm}$ are then orthogonal to all of these,
and hence to all $t_a v$, so these eigenvectors and the corresponding
eigenvalues are just the same as for the matrix
$\left.(\partial^2P(\phi)/\partial\phi_n\partial\phi_m)\right|_{\phi=v}$.
We see that for the unitarity gauge value $\xi\to\infty$, the Goldstone
bosons are so heavy that they drop out of the theory, while the other boson
masses are as usual.

The propagators are calculated by the usual rules: if the free-particle
part of the Lagrangian takes the form (after integration by parts)
$-\zeta^\dagger\mathcal D(\partial)\zeta$ for a complex field $\zeta$ or
$-\frac12\zeta^{\mathsf T}\mathcal D(\partial)\zeta$ for a real field $\zeta$, then
the propagator of this field is $\mathcal D^{-1}(ik)$. This gives the
propagators:
\begin{equation}
 A:\quad
 \Delta_{a\mu,b\nu}(k)
 =\left[
 \frac1{k^2+\mu^2}
 \left(\eta_{\mu\nu}
 -\frac{(1-\xi)k_\mu k_\nu}{k^2+\mu^2\xi}\right)
 \right]_{ab},
 \tag{21.2.20}\label{eq:21.2.20}
\end{equation}
\begin{equation}
 \phi:\quad
 \Delta_{nm}(k)
 =(k^2+M^2)^{-1}_{nm}
 +\xi\sum_{a,b}
 (t_a v)_n(t_b v)_m
 (k^2)^{-1}(k^2+\xi\mu^2)^{-1}_{ab},
 \tag{21.2.21}\label{eq:21.2.21}
\end{equation}
\begin{equation}
 \psi:\quad
 \Delta(k)=\frac{-i\sl{k}+m}{k^2+m^2},
 \tag{21.2.22}\label{eq:21.2.22}
\end{equation}
\begin{equation}
 \omega:\quad
 \Delta_{ab}(k)
 =(k^2+\xi\mu^2)^{-1}_{ab}.
 \tag{21.2.23}\label{eq:21.2.23}
\end{equation}
The poles in Eq.~\eqref{eq:21.2.20} at unphysical mass squares proportional
to $\xi$ are cancelled by the poles at the same masses in
Eq.~\eqref{eq:21.2.21}. Note that now for finite $\xi$ all propagators have
the same asymptotic behaviors as in the unbroken symmetry case, as required
for renormalizability. In particular, the $k_\mu k_\nu$ term in the vector
boson propagator no longer presents any problem for renormalizability,
because it is accompanied with an extra factor
$(k^2+\mu^2\xi)^{-1}$. We can even drop this term by choosing Feynman
gauge, with $\xi=1$. It is only in the unitarity gauge case where
$\xi\to\infty$ that this factor fails to give the propagator the asymptotic
behavior needed for renormalizability.

Even with well-behaved propagators, it is still necessary to show that the
ultraviolet divergences in these theories are constrained by the broken
gauge symmetries in such a way that every infinity can be cancelled by the
renormalization of a field or a parameter in the Lagrangian. This can be
done by the same techniques as in Sections 17.2 and 17.3, but treating the
vacuum expectation values of the spinless fields as external fields rather
than fixed quantities that break the
symmetries.~\hyperref[ch21-ref-5a]{[5a]}
